\def\stat{adam+volkov}





\def\tit{НЕСБАЛАНСИРОВАННОСТЬ КЛАССОВ\\ В~ТЕХНОЛОГИИ ПОДДЕРЖКИ\\  


КОНКРЕТНО-ИСТОРИЧЕСКИХ ИССЛЕДОВАНИЙ}





\def\titkol{Несбалансированность классов в~технологии ПКИИ} %поддержки  конкретно-исторических исследований}





\def\aut{И.\,М.~Адамович$^1$, О.\,И.~Волков$^2$}





\def\autkol{И.\,М.~Адамович, О.\,И.~Волков}





\titel{\tit}{\aut}{\autkol}{\titkol}





\index{Адамович И.\,М.}


\index{Волков О.\,И.}


\index{Adamovich I.\,M.}


\index{Volkov O.\,I.}





%{\renewcommand{\thefootnote}{\fnsymbol{footnote}}


%\footnotetext[1]


%{Исследование выполнено с~использованием ЦКП <<Информатика>> ФИЦ ИУ РАН.}} 





\renewcommand{\thefootnote}{\arabic{footnote}}


\footnotetext[1]{Федеральный исследовательский центр <<Информатика и~управление>> Российской академии наук, 


\mbox{Adam@amsd.com}}


\footnotetext[2]{Федеральный исследовательский центр <<Информатика и~управление>> Российской академии наук, 


\mbox{Volkov@amsd.com}}








 \label{st\stat}





 \thispagestyle{headings}


 


\vspace*{-14pt}








  





  \Abst{Статья продолжает серию работ, посвященных технологии поддержки  


кон\-крет\-но-ис\-то\-ри\-че\-ских исследований (ПКИИ), построенной на принципах сотворчества 


и~краудсорсинга и~ориентированной на широкий круг не относящихся 


к~профессиональным историкам и~биографам пользователей. Статья посвящена 


дальнейшему развитию темы подготовки данных для применяемых в~технологии 


алгоритмов машинного обучения. Показана особая важность бинарной классификации для 


кон\-крет\-но-ис\-то\-ри\-че\-ско\-го исследования. Описана проблема дисбаланса классов при 


проведении бинарной классификации с~по\-мощью алгоритмов машинного обучения и~ее 


последствия. Показано, что кон\-крет\-но-ис\-то\-ри\-че\-ские данные могут быть сильно 


несбалансированными. Приведен обзор подходов к~решению задачи устранения 


дисбаланса классов. Проведен анализ специфики конкретно-исторических данных, на его 


основе выбран подход oversampling как наиболее подходящий для технологии. Описаны 


алгоритмы, реализующие данный подход, оценены их достоинства и~недостатки. Выбран 


алгоритм ADASYN как наиболее перспективный для использования в~условиях 


технологии. Оценены возможности уже включенных в~технологию средств борьбы 


с~шумами в~данных и~выбросами для компенсации такого недостатка алгоритма 


ADASYN, как чувствительность к~выбросам.}


  


  \KW{конкретно-историческое исследование; распределенная технология; машинное 


обучение; дисбаланс классов; алгоритм ADASYN}





\DOI{10.14357\!/\!08696527230414}{YDVCYC}  





\vskip 10pt plus 9pt minus 6pt





 \thispagestyle{headings}


 


 \vspace*{-18pt}


  


\section{Введение}





\vspace*{-4pt}





  Поддержка конкретно-исторических исследований стала одной из 


актуальных задач современности, что обусловлено вовлечением 


в~исследовательский процесс не только членов профессионального 


исторического сообщества, но и~самых широких слоев непрофессионалов 


в~связи со все возрастающим интересом к~частной, семейной истории~[1].


  


  В~[2, 3] описана разработанная в~ФИЦ ИУ РАН распределенная 


технология ПКИИ, основанная на принципах краудсорсинга (мобилизации 


ресурсов ши-\linebreak\vspace*{-12pt}





\pagebreak





\noindent


рокого круга добровольцев посредством информационных 


технологий). Данные в~этой технологии организованы в~форме 


семантической сети. Узлы сети представляют собой именованные 


универсальные классы объектов. Факты задаются значениями экземпляров 


классов и~связями между ними. Связи наследуются из сети классов~[4]. 


  


  В конкретно-исторических исследованиях применяются следующие 


методы и~приемы работы с~данными~[5]:


  \begin{itemize}


  \item извлечение и~группировка изучаемых объектов; сравнение групп 


объектов; получение сводных данных об изучаемых группах объектов; 


\item классификация изучаемых объектов; 


  \item формирование хронологических рядов; периодизация (деление 


процессов развития на отличающиеся друг от друга периоды).


  \end{itemize}


  


  В~[6] показано, что подобные задачи целесообразно решать в~рамках 


технологии ПКИИ, в~том числе с~применением методов машинного 


обучения. Особое внимание обращено на процедуру очистки данных (data 


cleansing, data cleaning, data scrubbing или data wrangling), к~которой 


относятся~[7]:


  \begin{itemize}


  \item заполнение пропущенных значений;


  \item сглаживание зашумленных данных;


  \item устранение несогласованности;


  \item удаление выбросов.


  \end{itemize}


  


В~[6] также показана особая важность подготовки данных для задач 


машинного обучения в~рамках технологии ПКИИ в~связи 


с~фрагментарностью и~противоречивостью кон\-крет\-но-ис\-то\-ри\-че\-ской 


информации. Описана специфика очистки кон\-крет\-но-ис\-то\-ри\-че\-ских 


данных и~проведен анализ возможности применения с~этой \mbox{целью} 


механизмов и~алгоритмов, уже интегрированных в~технологию. Перечислены 


основные направления, по которым проводится очистка данных. Для каждого 


направления выявлены подходящие инструменты, уже включенные 


в~технологию.


  


  Данная статья посвящена дальнейшему развитию технологии ПКИИ 


в~об\-ласти подготовки данных для задач, связанных с~машинным  


обуче\-нием.





\section{Дисбаланс классов}





  В~[8] показана значимость классификации и~типологизации, которая по 


форме представляет собой разновидность классификации, лежащих в~основе  


ис\-то\-ри\-ко-ти\-по\-ло\-ги\-че\-ско\-го метода исторического познания. 


  


  Особую важность в~конкретно-историческом исследовании имеет 


бинарная классификация, поскольку она соответствует подтверждению или 


отбрасыванию гипотезы, возникшей в~процессе исследования относительно 


объекта исследования. Например, если историческое исследование 


посвящено отношению российских дворян к~советской власти, то одно из 


возможных дихотомических делений в~соответствии с~[9] может выглядеть 


следующим образом:





\smallskip


    


    Русские дворяне с~1917~г.:


    \begin{enumerate}[1.]


    \item Приняли светскую власть.


    \begin{enumerate}[{1}.1.]


    \item Участвовали в~гражданской войне.


    \begin{enumerate}[{1.1.}1.]


    \item Погибли.


    \item Не погибли.


    \end{enumerate}


    \item Не участвовали в~гражданской войне.


    \end{enumerate}


    \item Не приняли светскую власть.


    \begin{enumerate}[{2}.1.]


    \item Участвовали в~гражданской войне.


    \begin{enumerate}[{2.1.}1.]


    \item  Погибли.


    \item Не погибли.


    \begin{enumerate}[{2.1.2.}1.]


    \item Эмигрировали.


    \item Остались.


    \end{enumerate}


    \end{enumerate}


    \item Не участвовали в~гражданской войне.


    \begin{enumerate}[{2.2.}1.]


    \item  Эмигрировали.


        \item Остались.


    \end{enumerate}


    \end{enumerate}


    \end{enumerate}


    


  Но в~задачах классификации данным присущи дополнительные  


проб\-ле\-мы помимо описанных выше проблем очистки данных. В~задаче 


бинарной классификации данные называются несбалансированными 


(Imbalanced Data), если большинство объектов исходного набора данных 


принадлежат одному классу (мажоритарный класс) и~гораздо меньшее число 


объектов относится к~другому классу (миноритарный класс). При этом 


неверная классификация примеров миноритарного класса, как правило, 


обходится в~разы дороже, чем ошибочная классификация примера 


мажоритарного класса, так как экземпляры миноритарного класса 


представляют собой редкие, но наиболее важные данные в~реально 


используемых наборах данных~[10]. Несбалансированность классов 


считается сильной, когда классы в~выборке представлены в~соотношении 


более чем $10:1$. При обуче\-нии модели на сильно несбалансированных 


данных алгоритм классификации может посчитать примеры интересующего 


миноритарного класса шумом и~пропустить их~[11].


  


  Специфика конкретно-исторических данных состоит в~том, что они 


охватывают самые разные аспекты жизни и~обладают большой 


неоднородностью. Например, согласно переписи 1897~г., все население 


страны, а~это 125~млн человек, распределялось на следующие 


сословия: дворяне~--- 1,5\% всего населения; духовенство~--- 0,5\%;  


куп\-цы~--- 0,3\%; мещане~--- 10,6\%; крестьяне~--- 77,1\%; казаки~--- 2,3\%. 


А~в~годы революции и~гражданской войны интеллигенция составляла 


около~2,5\% по отношению ко всему населению~[12]. Как видно из 


приведенных примеров, проб\-ле\-ма сильного дисбаланса классов в~полной 


мере присуща кон\-крет\-но-ис\-то\-ри\-че\-ским данным и,~используя 


методы машинного обучения при решении задач классификации, следует 


принимать меры по ее решению.


  


\section{Подходы к~устранению дисбаланса классов}





  Можно выделить три основных подхода к~решению указанной  


проб\-ле\-мы~[11]. 


  \begin{enumerate}[1.]


  \item Первый подход связан с~изменением распределения классов 


в~выборке. Эта задача решается с~помощью определенного отбора, т.\,е.\ 


семплинга (sampling), примеров из выборки таким образом, чтобы 


соотношение числа примеров в~классах соответствовало заданным 


требованиям. 


  \item Ко второму подходу относятся все методы распределения весов. Этот 


принцип основан на присвоении каждому классу числового параметра~--- 


веса, который соответствует штрафу за неверную классификацию. Вес 


определяется экспертной оценкой либо исходя из распределения классов 


в~обуча\-ющей выборке. Этот подход нельзя считать абсолютно 


универсальным, поскольку он опирается на поддержку весов классов 


алгоритмами. Например, популярная библиотека sklearn предусматривает 


аргумент class\_weight только для алгоритмов решающего дерева, случайного 


леса и~логистической регрессии.


\item В третьем подходе примеры миноритарного класса считаются 


исключительными и~классификация строится в~рамках задачи обнаружения 


аномалий. Этот подход применяется, когда объектов миноритарного класса 


в~обуче\-нии крайне мало, например~1--3 на 1~млн объектов (а~такие 


задачи бывают, например, когда позитивный класс~--- катастрофы или 


большие экономические кризисы)~[13], что не соответствует специфике  


кон\-крет\-но-ис\-то\-ри\-че\-ских данных, как видно из приведенных выше 


примеров.


  \end{enumerate}


  


  На практике часто используется идея семплинга (первый подход) ввиду ее 


простоты и~разнообразия. В~этом случае применяют либо добавление 


в~обуча\-ющую выборку примеров миноритарного класса (oversampling), 


либо удаление из нее примеров мажоритарного класса (undersampling) 


с~\mbox{целью} достичь сбалансированного распределения.


  


\section{Алгоритмы генерации примеров}





  Поскольку конкретно-исторические данные обладают свойством 


фрагментарности, а~их поиск становится наиболее трудоемким этапом 


исследования~[14], концепция удаления данных не соответствует специфике 


кон\-крет\-но-ис\-то\-ри\-че\-ско\-го исследования, и~для целей технологии ПКИИ 


следует рассматривать варианты, связанные с~генерацией примеров 


миноритарного класса (oversampling). Поскольку применение oversampling 


может привести к~переобучению~[15], при выборе конкретного алгоритма 


машинного обучения потребуются дополнительные меры по устранению 


данного недостатка с~учетом специфики кон\-крет\-но-ис\-то\-ри\-че\-ских 


данных. В~част\-ности, из принципиальной ограниченности объема этих 


данных вытекает необходимость использования простых моделей, поскольку 


сложные модели более склонны к~переобучению на ограниченных выборках. 


И~для этих моделей потребуется поиск баланса предикторов и~размера 


выборки, поскольку переобучение вероятно даже для простых моделей, если 


подать на вход модели избыточное число предикторов.


  


  К наиболее популярным алгоритмам oversampling относятся~[16]:


  \begin{enumerate}[(1)]


  \item случайная наивная выборка.


  


  Самый простой способ увеличить чис\-ло примеров миноритарного 


класса~--- случайным образом выбрать наблюдения из него и~добавить их 


в~общий набор данных, пока не будет достигнут необходимый баланс между 


миноритарным и~мажоритарным классами. К~достоинствам такого подхода 


относятся его простота, легкость реализации и~предоставляемая им 


возможность изменить баланс в~любую нуж\-ную сторону. Такой подход 


к~восстановлению баланса не всегда эффективен, поэтому был предложен 


следующий метод увеличения числа примеров миноритарного класса;


  \item алгоритм SMOTE. 


  


  Алгоритм SMOTE (Synthetic Minority Over-sampling Technique)~[17] 


синтезирует элементы миноритарного класса в~непосредственной бли\-зости 


от уже существующих элементов. Алгоритм находит разность между данным 


элементом~$x_i$ и~его ближайшим соседом~$\hat{x}_i$, выбранным 


случайно из группы~$K$~ближайших соседей, найденным с~по\-мощью 


алгоритма k-NN (k-nearest neighbors algorithm). Эта разность умножается на 


случайное число~$\delta$ в~интервале от~0 до~1. Полученное значение 


добавляется к~данному образцу для формирования нового синтезированного 


образца~$x_{\mathrm{new}}$ в~пространстве признаков:


  $$


  x_{\mathrm{new}} =x_i+\left( \hat{x}_i -x_i\right) \delta\,.


  $$


  


  Подобные действия продолжаются со следующим соседом до заданного 


пользователем числа образцов. Существует ряд модификаций этого 


алгоритма. 


  


  Алгоритм SMOTE решает многие проблемы, которые присущи методу 


случайной наивной выборки, но имеет и~свои недостатки, главный из 


которых~--- игнорирование мажоритарного класса. Это может проявиться 


в~том, что при сильно разреженном распределении объектов миноритарного 


класса относительно мажоритарного наборы данных <<смешаются>>, т.\,е.\ 


расположатся в~таком виде, что отделить объекты одного класса от другого 


станет очень трудно. Такая ситуация в~условиях технологии ПКИИ вполне 


возможна из за принципиальной фрагментарности  


кон\-крет\-но-ис\-то\-ри\-че\-ских данных. В~качестве решения этой  


проб\-ле\-мы был предложен третий метод;


  \item алгоритм ADASYN.


  


  Основная особенность алгоритма ADASYN (Adaptive Synthetic 


Sampling)~[18] заключается в~том, что он использует специальный механизм 


для автоматического определения числа синтетических образцов, которые 


должен сгенерировать каждый элемент миноритарного класса обучающей 


выборки, вместо того чтобы синтезировать одинаковое число образцов для 


каждого элемента миноритарного класса, как SMOTE. Пусть имеется 


выборка, содержащая~$m$~наблюдений~$x_i$, $i\hm=1,\ldots , m$. Алгоритм 


выполняется в~соответствии со следующими шагами:


  \begin{itemize}


  \item[(а)] для заданной пропорции классов вычисляется общее число 


синтези\-ру\-емых элементов


  $$


  G=\left( S_{\mathrm{maj}}-S_{\mathrm{min}}\right) \beta\,,


  $$


где $S_{\mathrm{maj}}$ и~$S_{\mathrm{min}}$~--- число наблюдений в~мажоритарном 


и~миноритарном классах соответственно; $\beta\hm \in [0,1]$~--- уровень 


баланса выборки. Если $\beta\hm=1$, то после синтеза выборка станет 


полностью сбалансированной;


  \item[(б)] для каждого наблюдения~$x_i$ из миноритарного класса  


с~по\-мощью алгоритма k-NN выявляются ее $K$~ближайших соседей 


и~вычисляется отношение $r_i\hm= \Delta_i/K$, где $\Delta_i$~--- число 


элементов мажоритарного класса среди ближайших $K$~соседей~$x_i$, т.\,е.\ 


$r_i\hm\in [0,1]$;


  \item[(в)]  проводится нормирование~$r_i$ в~соответствии с~выражением


  $$


  \hat{r}_i=\fr{r_i}{\sum\nolimits^m_{j=1} r_j}\,;


  $$


  \item[(г)] определяется число $g_i$ искусственных элементов, которые должны 


быть синтезированы для каждого наблюдения~$x_i$ из миноритарного 


класса


  $$


  g_i=\hat{r}_i G\,;


  $$


  \item[(д)] проводится генерация в~соответствии с~алгоритмом SMOTE.


  \end{itemize}


  


  Основная идея алгоритма ADASYN состоит в~том, что он 


использует~$\hat{r}_i$ для определения числа элементов, которые должны 


быть синтезированы для каждого элемента миноритарного класса. Это 


эквивалентно присвоению весовых коэффициентов. Чем больше 


окружающих элементов мажоритарного класса, тем выше вес.


  


  Недостаток ADASYN заключается в~том, что на него легко влияют 


выбросы. Если все $K$ ближайших соседей элемента миноритарного класса 


окажутся элементами мажоритарного класса, его вес станет довольно 


большим и~вокруг него будет сгенерировано больше искусственных 


наблюдений~[19]. Из принципиальной неоднородности кон\-крет\-но-ис\-то\-ри\-че\-ских 


данных следует, что такая ситуация в~условиях технологии 


ПКИИ вполне возможна. При этом выбросы могут быть как ошибочными 


данными в~связи с~оби\-ли\-ем недостоверной и~искаженной информации, так 


и~реальными данными с~крайне редкими сочетаниями признаков. В~[6] 


показано, что технология ПКИИ уже имеет в~своем составе средства борьбы с~шумами в~данных и~выбросами, включенные в~технологию в~рамках 


механизма выявления аномалий в~конкретно-исторических данных~[20]. 


Поэтому алгоритм ADASYN представляется более перспективным для 


использования в~технологии.


  \end{enumerate}


  


\section{Выводы}





  Как показано выше, методы классификации и~типологизации активно 


применяются в~историческом исследовании. Особую значимость  


в~кон\-крет\-но-ис\-то\-ри\-че\-ском исследовании имеет бинарная 


классификация. Ее выполнение в~технологии ПКИИ может осуществляться 


с~применением алгоритмов машинного обучения. Этап подготовки данных 


для этих алгоритмов помимо очистки данных должен включать устранение 


дисбаланса классов, поскольку проблема сильного дисбаланса классов 


в~полной мере присуща конкретно-историческим данным, а~при обуче\-нии 


модели на сильно несбалансированных данных алгоритм классификации 


может посчитать примеры интересующего миноритарного класса шумом 


и~пропустить их. Анализ методов устранения дисбаланса классов в~условиях 


технологии ПКИИ показал потенциальную эффективность для этой цели 


алгоритма ADASYN в~сочетании с~включенными в~технологию средствами 


борьбы с~шумами в~данных и~выбросами.


  


  Актуальность дальнейшего развития технологии ПКИИ, ориентированной на 


широкий круг не относящихся к~профессиональным историкам  


и~био\-гра\-фам пользователей, вытекает из все воз\-рас\-та\-юще\-го 


общественного интереса к~частной, семейной истории.


  


{\small\frenchspacing


{%\baselineskip=10.8pt


\addcontentsline{toc}{section}{Литература}


\begin{thebibliography}{99}


  \bibitem{1-ad}


  \Au{Грибач С.\,В.} Исследование семейных кризисов посредством 


психолингвистического эксперимента~/\!\!/ Сборники конф. НИЦ Социосфера, 2010. №\,6. 


С.~45--54.


  \bibitem{2-ad}


  \Au{Адамович И.\,М., Волков О.\,И.} Технология распределенного автоматизированного 


анализа исторических текстов~/\!\!/ Системы и~средства информатики, 2016. Т.~26. №\,3. 


С.~148--161. doi: 10.14357\!/\!08696527160311.


  \bibitem{3-ad}


  \Au{Адамович И.\,М., Волков О.\,И.} Единая технология поддержки  


кон\-крет\-но-ис\-то\-ри\-че\-ских исследований~/\!\!/ Системы и~средства информатики, 


2019. Т.~29. №\,1. С.~194--205. doi: 10.14357\!/\!08696527190116.


  \bibitem{4-ad}


  \Au{Адамович И.\,М., Волков О.\,И.} Принципы организации данных для технологии 


поддержки кон\-крет\-но-ис\-то\-ри\-че\-ских исследований~/\!\!/ Системы и~средства 


информатики, 2019. Т.~29. №\,2. С.~161--171. doi: 10.14357\!/\!08696527190214.


  \bibitem{5-ad}


  \Au{Гагарина Д.\,А., Корниенко С.\,И., Поврозник~Н.\,Г.} Информационные системы 


в~циф\-ро\-вой среде исторической науки~/\!\!/ История, 2017. Т.~7. Вып.~7(51). 12~с.


  \bibitem{6-ad}


  \Au{Адамович И.\,М., Волков О.\,И.} Очистка данных в~технологии поддержки 


кон\-крет\-но-ис\-то\-ри\-че\-ских исследований~/\!\!/ Системы и~средства информатики, 2023. 


Т.~33. №\,3. С.~149--160. doi: 10.14357\!/\!08696527230313.


  \bibitem{7-ad}


  \Au{Osborne J.\,W.} Best practices in data cleaning: A~complete guide to everything you 


need to do before and after collecting your data.~--- Newbury Park, CA, USA: SAGE 


Publications Inc., 2012. 275~p.


  \bibitem{8-ad}


  \Au{Мартюшов Л.\,Н.} Методы исторического исследования.~--- 


Екатеринбург: УрГПУ, 2016. 86~с.


  \bibitem{9-ad}


  \Au{Бочаров А.\,В.} Алгоритмы использования основных научных методов  


в~кон\-крет\-но-ис\-то\-ри\-че\-ском исследовании.~--- Томск: ТГУ, 


2007. 137~с.


  \bibitem{10-ad}


  \Au{Клюева И.\,А.} Исследование применимости SMOTE-ал\-го\-рит\-ма при 


классификации несбалансированных данных~/\!\!/ Современные технологии в~науке 


и~образовании: Сб. трудов II~Междунар. научн.-тех\-нич. и~научн.-ме\-то\-ди\-че\-ской конф.~/ Под ред.\ О.\,В.~Миловозова.~--- Рязань: РГРТУ, 2017. Т.~1. С.~143--146.


  EDN: ZCKRHN.


  \bibitem{11-ad}


  \Au{Найденов А.\,С.} Применение сэмплинга в~условиях несбалансированности 


классов~/\!\!/ Информационные технологии: Межвузовский сборник научных трудов.~--- 


Рязань: РГРТУ, 2017. С.~73--75. EDN: ZQUNVJ.


  \bibitem{12-ad}


  \Au{Фёдоров А.\,Н.} Реальная опора советской власти:  


со\-ци\-аль\-но-де\-мо\-гра\-фи\-че\-ские характеристики городского населения России 


в~1917--1920~годах (на материалах Центрального промышленного района)~/\!\!/ 


Ж.~исследований социальной политики, 2010. Т.~8. №\,1. С.~69--86. EDN: LDFPAF.


  \bibitem{13-ad}


  \Au{Дьяконов А.\,Г.} Дисбаланс классов~/\!\!/ Анализ малых данных, 2021. {\sf 


https://\linebreak alexanderdyakonov.wordpress.com/2021/05/27/imbalance}.


  \bibitem{14-ad}


  \Au{Адамович И.\,М., Волков О.\,И.} Иерархическая форма представления 


биографического факта~/\!\!/ Системы и~средства информатики, 2016. Т.~26. №\,2.  


С.~108--122. doi: 10.14357\!/\!08696527160207.


  \bibitem{15-ad}


  \Au{Japkowicz N., Stephen~S.} The class imbalance problem: A~systematic study~/\!\!/ 


Intell. Data Anal., 2002. Vol.~6. Iss.~5. P.~429--449. 


  \bibitem{16-ad}


  \Au{Севастьянов Л.\,А., Щетинин Е.\,Ю.} О~методах повышения точ\-ности 


многоклассовой классификации на несбалансированных данных~/\!\!/ Информатика и~её 


применения, 2020. Т.~14. Вып.~1. С.~63--70. doi: 10.14357\!/\!19922264200109.


  \bibitem{17-ad}


  \Au{Chawla N.\,V., Bowyer K.\,W., Hall L.\,O., Kegelmeyer~W.\,P.} SMOTE: Synthetic 


minority over-sampling technique~/\!\!/ J.~Artif. Intell. Res., 2002. Vol.~16.  


P.~321--357. doi: 10.1613\!/\!jair.953.


  


  \bibitem{18-ad}


  \Au{He H., Bai Ya., Garcia~E.\,A., Li~Sh.} ADASYN: Adaptive synthetic sampling approach 


for imbalanced learning~/\!\!/ IEEE Joint Conference (International) on Neural Networks (IEEE 


World Congress on Computational Intelligence).~--- Piscataway, NJ, USA: IEEE, 2008. 


P.~1322--1328. doi: 10.1109\!/\!IJCNN.2008.4633969.


  \bibitem{19-ad}


  Category imbalance problems of machine learning~(3)~--- sampling method~/\!\!/ 


Cnblogs.com: Massquantity blog, 2018. {\sf https://www.cnblogs.com/massquantity/\linebreak p/9382710.html}.


  \bibitem{20-ad}


  \Au{Адамович И.\,М., Волков О.\,И.} Алгоритмы кластеризации для технологии 


поддержки кон\-крет\-но-ис\-то\-ри\-че\-ских исследований~/\!\!/ Системы и~средства 


информатики, 2022. Т.~32. №\,4. С. 112--123. doi: 10.14357\!/\!08696527220411.





   \end{thebibliography}


} }











\vspace*{-6pt}





\hfill{\small\textit{Поступила в~редакцию 20.07.23}}





\vspace*{8pt}





\hrule





\vspace*{2pt}











\hrule





\vspace*{8pt}





%\newpage





%\vspace*{-28pt}





\def\tit{CLASS IMBALANCE IN~THE~TECHNOLOGY\\ OF~CONCRETE HISTORICAL 


INVESTIGATION SUPPORT}











\def\titkol{Class imbalance in~the~technology of~concrete historical 


investigation support}





\def\aut{I.\,M.~Adamovich and O.\,I.~Volkov}





\def\autkol{I.\,M.~Adamovich and O.\,I.~Volkov}





\titel{\tit}{\aut}{\autkol}{\titkol}





\vspace*{-8pt}





\noindent


Federal Research Center ``Computer Science and Control'' of the Russian Academy


of Sciences,   44-2~Vavilov Str., Moscow 119133, Russian Federation








\def\leftfootline{\small{\textbf{\thepage}


\hfill Sistemy i Sredstva Informatiki~---


Systems and Means of Informatics 2023 vol~33 no\,4}


}%


 \def\rightfootline{\small{Sistemy i Sredstva Informatiki~---


 Systems and Means of Informatics 2023 vol~33 no\,4


\hfill \textbf{\thepage}}}





\vspace*{3pt}


  


  





\Abste{The article continues a series of works devoted to the technology of concrete historical 


investigation support, built on the principles of co-creation and crowdsourcing and designed for 


a~wide range of nonprofessional historians and biographers users. The article is 


devoted to the further development of the topic of data preparation for machine learning 


algorithms used in the technology. The special importance of binary classification for concrete 


historical research is shown. The problem of class imbalance in binary 


classification using machine learning algorithms and its consequences are described. It 


is shown that concrete historical data can be highly imbalanced. An overview of approaches to 


solving the problem of class imbalance elimination is given. The analysis of the specifics of concrete 


historical data was carried out, and on its basis, the oversampling approach was chosen as the 


most suitable for the technology. Algorithms implementing this approach are described; their 


advantages and disadvantages are evaluated. The ADASYN algorithm has been selected as the 


most promising for use in the technology conditions. The possibilities of the already included in the technology


means of data noise and 


outliers control  to compensate such a~disadvantage of the 


ADASYN algorithm as sensitivity to outliers are evaluated.}





\KWE{concrete historical investigation; distributed technology; machine learning; class 


imbalance; ADASYN algorithm}





\DOI{10.14357\!/\!08696527230414}{YDVCYC}  





%\vspace*{-10pt}








%\Ack


%\noindent





%\vspace*{-3pt}





\renewcommand{\bibname}{\protect\rmfamily References}


%\renewcommand{\bibname}{\large\protect\rm References}





{\small\frenchspacing


{ %\baselineskip=12pt


\addcontentsline{toc}{section}{References}


\begin{thebibliography}{99}


  \bibitem{1-ad-1}


  \Aue{Gribach, S.\,V.} 2010. Is\-sle\-do\-va\-nie se\-mey\-nykh kri\-zi\-sov po\-sred\-st\-vom 


psi\-kho\-ling\-vi\-sti\-che\-sko\-go eks\-pe\-ri\-men\-ta [The study of family crises through a psycholinguistic 


experiment]. \textit{Sborniki konf. NITs Sotsiosfera} [Conference Proceedings NIC Sociosfera] 


6:45--54.


  \bibitem{2-ad-1}


  \Aue{Adamovich, I.\,M., and O.\,I.~Volkov.} 2016. Tekh\-no\-lo\-giya ras\-pre\-de\-len\-no\-go 


av\-to\-ma\-ti\-zi\-ro\-van\-no\-go ana\-li\-za is\-to\-ri\-che\-skikh teks\-tov [The distributed automated technology of 


historical texts analysis]. \textit{Sistemy i~Sredstva Informatiki~--- Systems and Means of 


Informatics} 26(3):148--161. doi: 10.14357\!/\!08696527160311.


  \bibitem{3-ad-1}


  \Aue{Adamovich, I.\,M., and O.\,I.~Volkov.} 2019. Edi\-naya tekh\-no\-lo\-giya pod\-derz\-hki 


konkretno-istoricheskikh is\-sle\-do\-va\-niy [Unified technology of concrete historical research 


support]. \textit{Sistemy i~Sredstva Informatiki~--- Systems and Means of Informatics} 


29(1):194--205. doi: 10.14357\!/\!08696527190116.


  \bibitem{4-ad-1}


  \Aue{Adamovich, I.\,M., and O.\,I.~Volkov.} 2019. Prin\-tsi\-py or\-ga\-ni\-za\-tsii dan\-nykh dlya 


tekh\-no\-lo\-gii pod\-derz\-hki konkretno-istoricheskikh is\-sle\-do\-va\-niy [The principles of data 


organization for the technology of concrete historical research support]. \textit{Sistemy 


i~Sredstva Informatiki~--- Systems and Means of Informatics} 29(2):161--171. doi: 


10.14357\!/\!08696527190214.


  \bibitem{5-ad-1}


  \Aue{Gagarina, D.\,A., S.\,I.~Korniyenko, and N.\,G.~Povroznik.} 2017. In\-for\-ma\-tsi\-on\-nye 


sis\-te\-my v~tsif\-ro\-voy sre\-de is\-to\-ri\-che\-skoy na\-u\-ki [Information systems in the digital environment of 


historical studies]. \textit{Istoriya} [Hystory] 


7(51). 12~p. Available at: {\sf https://history.jes.su/s207987840001638-0-1/} (accessed October~16, 2023).


  \bibitem{6-ad-1}


  \Aue{Adamovich, I.\,M., and O.\,I.~Volkov.} 2023. Ochistka dannykh v tekhnologii 


podderzhki konkretno-istoricheskikh issledovaniy [Data cleansing in the technology of concrete 


historical investigation support]. \textit{Sistemy i~Sredstva Informatiki~--- Systems and Means 


of Informatics} 33(3):149--160. doi: 10.14357\!/\!08696527230313.


  \bibitem{7-ad-1}


  \Aue{Osborne, J.\,W.} 2012. \textit{Best practices in data cleaning: A~complete guide to 


everything you need to do before and after collecting your data}. Newbury Park, CA: SAGE 


Publications Inc. 275~p.


  \bibitem{8-ad-1}


  \Aue{Martyushov, L.\,N.} 2016. \textit{Me\-to\-dy is\-to\-ri\-che\-sko\-go is\-sle\-do\-va\-niya} 


  [Methods of historical research]. Ekaterinburg: USPU. 86~p.


  \bibitem{9-ad-1}


  \Aue{Bocharov, A.\,V.} 2007. \textit{Al\-go\-rit\-my is\-pol'\-zo\-va\-niya osnov\-nykh na\-uch\-nykh 


me\-to\-dov v~konkretno-istoricheskom is\-sle\-do\-va\-nii} [Algorithms for the use of basic scientific methods 


in concrete historical research]. Tomsk: Publishing House of the Tomsk State University. 140~p.


  \bibitem{10-ad-1}


  \Aue{Klyuyeva, I.\,A.} 2017. Is\-sle\-do\-va\-nie pri\-me\-ni\-mosti SMOTE-algoritma pri klas\-si\-fi\-ka\-tsii 


nesba\-lan\-si\-ro\-van\-nykh dan\-nykh [The study of the applicability of the SMOTE-algorithm by 


classification of the imbalanced data]. \textit{So\-vre\-men\-nye tekh\-no\-lo\-gii v~nauke i~ob\-ra\-zo\-va\-nii: 


Sbornik tr. II Mezhdunar. nauchn.-tekhnich. i~nauchn.-metodicheskoy konf.} [Modern 


Technologies in Science and Education: 2nd  Scientific-Technical and 


Scientific-Methodological Conference (International) Proceedings]. Ed.\ O.\,V.~Milovzorov. Ryazan: Ryazan State 


Radioengineering University.  1:143--146. EDN: ZCKRHN.


  \bibitem{11-ad-1}


  \Aue{Naydenov, A.\,S.} 2017. Pri\-me\-ne\-nie semp\-lin\-ga v~uslo\-vi\-yakh nesba\-lan\-si\-ro\-van\-nosti 


klas\-sov [Application of sampling in conditions of unbalanced classes]. \textit{Informatsionnye 


tekhnologii: Mezhvuzovskiy sbornik nauchnykh trudov} [Information technologies: Interuniversity 


collection of scientific papers]. Ryazan: Ryazan State Radioengineering University. 1:73--75. EDN: 


ZQUNVJ.


  \bibitem{12-ad-1}


  \Aue{Fedorov, A.\,N.} 2010. Re\-al'\-naya opo\-ra so\-vets\-koy vlas\-ti: sotsial'no-demograficheskie 


kha\-rak\-te\-ri\-sti\-ki go\-rod\-sko\-go na\-se\-le\-niya Ros\-sii v~1917--1920 go\-dakh (na materialakh 


Tsentral'nogo promyshlennogo rayona) [The real support of Soviet power: Socio-demographic 


characteristics of the urban population of Russia in 1917--1920 (on the materials of the Central 


Industrial District)]. \textit{Zh.~is\-sle\-do\-va\-niy so\-tsi\-al'\-noy po\-li\-ti\-ki} [J.~Social Policy 


Studies] 1(8):69--86. EDN: LDFPAF.


  \bibitem{13-ad-1}


  \Aue{Dyakonov, A.\,G.} 2021. Disbalans klas\-sov [Class imbalance]. \textit{Ana\-liz ma\-lykh 


dan\-nykh} [Small data analysis]. Available at: {\sf 


https://alexanderdyakonov.wordpress.com/\linebreak  2021/05/27/imbalance/} (accessed October~16, 2023).


  \bibitem{14-ad-1}


  \Aue{Adamovich, I.\,M., and O.\,I.~Volkov.} 2016. Ie\-rar\-khi\-che\-skaya for\-ma pred\-stav\-le\-niya 


bio\-gra\-fi\-che\-sko\-go fak\-ta [Hierarchial format of a~biographical fact]. \textit{Sis\-te\-my i~Sred\-st\-va 


Informatiki~--- Systems and Means of Informatics} 26(2):108--122. doi:  


10.14357\!/\!08696527160207.


  \bibitem{15-ad-1}


  \Aue{Japkowicz, N., and S.~Stephen.} 2002. The class imbalance problem: A~systematic 


study. \textit{Intell. Data Anal.} 6(5):429--449.


  \bibitem{16-ad-1}


  \Aue{Sevastianov, L.\,A., and E.\,Yu.~Shchetinin.} 2020. O~me\-to\-dakh po\-vy\-she\-niya 


toch\-nosti mno\-go\-klas\-so\-voy klas\-si\-fi\-ka\-tsii na nesba\-lan\-si\-ro\-van\-nykh dan\-nykh [On methods for 


improving the accuracy of multiclass classification on imbalanced data]. \textit{Informatika i~ee 


Primeneniya~--- Informatics and Applications} 14(1):63--70. doi: 


10.14357\!/ 19922264200109.


  \bibitem{17-ad-1}


  \Aue{Chawla, N.\,V., K.\,W. Bowyer, L.\,O.~Hall, and W.\,P.~Kegelmeyer.} 2002. SMOTE: 


Synthetic minority over-sampling technique. \textit{J.~Artif. Intell. Res.} 16:321--357. doi: 


10.1613\!/\!jair.953.


  \bibitem{18-ad-1}


  \Aue{He, H., Ya.~Bai, E.\,A.~Garcia, and Sh.~Li.} 2008. ADASYN: Adaptive synthetic 


sampling approach for imbalanced learning. \textit{IEEE Joint Conference (International) on Neural Networks (IEEE World Congress on 


Computational Intelligence)}. Piscataway, NJ: IEEE. 1322--1328. doi: 


10.1109\!/\!IJCNN.2008.4633969.


  \bibitem{19-ad-1}


  Category imbalance problems of machine learning (3)~--- sampling method. \textit{\mbox{Cnblogs.com}: Massquantity blog}. Available 


at: {\sf https://www.cnblogs.com/\linebreak massquantity/p/9382710.html} (accessed October~16, 2023).


  \bibitem{20-ad-1}


  \Aue{Adamovich, I.\,M., and O.\,I.~Volkov.} 2022. Al\-go\-rit\-my klas\-te\-ri\-za\-tsii dlya tekh\-no\-lo\-gii 


pod\-derzh\-ki konkretno-istoricheskikh is\-sle\-do\-va\-niy [Clustering algorithms for technology of concrete historical 


investigation support]. \textit{Sistemy i~Sredstva Informatiki~--- Systems and Means of 


Informatics} 32(4):112--123. doi: 10.14357\!/\!08696527220411.





\end{thebibliography}


} }





\vspace*{-6pt}





\hfill{\small\textit{Received July 20, 2023}} 





\vspace*{-16pt} 


  


\Contr





\vspace*{-3pt}





\noindent


\textbf{Adamovich Igor M.} (b.\ 1934)~--- Candidate of Science (PhD) in technology, leading 


scientist, Federal Research Center ``Computer Science and Control'' of the Russian Academy of 


Sciences, 44-2~Vavilov Str., Moscow 119133, Russian Federation; \mbox{adam@amsd.com}





\vspace*{3pt}





\noindent


\textbf{Volkov Oleg I.} (b.\ 1964)~--- leading programmer, Federal Research Center ``Computer 


Science and Control'' of the Russian Academy of Sciences, 44-2~Vavilov Str., Moscow 119133, 


Russian Federation; \mbox{volkov@amsd.com}





  


\label{end\stat}





\renewcommand{\bibname}{\protect\rm Литература}    


   


